\section{Por qué el sistema de coordenadas es el núcleo de la Part-Whole Relation}

Si un nostril respalda que «aquí hay una nose», no solo debe predecir la identity, sino también la position, la orientation y la scale de la nose. Dicho de otro modo, la part-whole relation no es una etiqueta discreta, sino una transformación entre intrinsic coordinate frames. Hinton recurre primero a un conjunto de experimentos psicológicos para que el público perciba que nuestra comprensión de un mismo conjunto de líneas cambia según el reference frame interno.

\lecturefigure{slide-08-psychological-coordinate-frames.jpg}{Las siguientes diapositivas usan demostraciones psicológicas para mostrar que el parse visual y los sistemas de coordenadas son reales.}{Video oficial de Stanford Online 00:09:16--00:13:31.}

\subsection{Cube demonstration: la dificultad no es la rotación, sino el cambio de frame}

En clase se pide al público imaginar un wire-frame cube, girar la body diagonal original hasta dejarla vertical y luego señalar la posición de las demás esquinas. Muchas personas no tienen dificultades con la geometría tridimensional; lo que ocurre es que están acostumbradas a ubicar el cube en un intrinsic frame de «cara inferior horizontal, aristas verticales hacia arriba». En cuanto se les obliga a usar un frame desconocido, la internal representation del mismo objeto cambia drásticamente.

\lecturefigure{slide-09-cube-demonstration.jpg}{Cube mental-rotation task: convertir la body diagonal conocida en el eje vertical.}{Video oficial de Stanford Online 00:13:32--00:16:07.}

\begin{knowledgebox}{Lectura de la figura: lo que la tarea cambia en realidad es el intrinsic frame}
Primero se localiza el front-bottom-right corner que originalmente toca la mesa y luego el top-back-left corner opuesto; en clase se pide girar esta body diagonal hasta convertirla en el eje vertical. Las relaciones geométricas de los otros seis corners no desaparecen, pero se retira el sistema de coordenadas conocido de «cara inferior horizontal, aristas verticales hacia arriba». Si al público le cuesta ubicar las esquinas, la evidencia respalda la reference-frame dependence, no una falta general de capacidad de rotación tridimensional.
\end{knowledgebox}

\begin{knowledgebox}{Qué pueden demostrar los experimentos psicológicos}
Muestran que el reasoning humano elige un reference frame y depende de él, y que cambiar de frame modifica las relaciones accesibles; no demuestran que las neuronas usen necesariamente GLOM embedding, matrix capsule ni alguna attention rule específica. Aquí el papel de la evidencia es de motivation, no de implementation identification.
\end{knowledgebox}

\subsection{Six rods: una misma entrada puede tener parses completamente distintos}

Las seis varillas pueden verse como una crown de tres lóbulos o como un central rectangle con varios lóbulos inclinados. Las dos percepts no discrepan sobre la posición de los segmentos externos, pero difieren por completo en qué aristas «se agrupan de forma natural» y cuáles parecen parallel. Hinton lo compara con la structural ambiguity de una oración: la truth condition puede ser la misma, pero el parse es distinto.

\lecturefigure{slide-10-six-rods-percept.jpg}{Otra percept de las seis varillas: una misma entrada geométrica puede organizarse con estructuras distintas.}{Video oficial de Stanford Online 00:16:08--00:16:55.}

\lecturefigure{slide-11-alternative-percepts.jpg}{Distintas internal representations pueden explicar un mismo hecho externo.}{Video oficial de Stanford Online 00:16:56--00:17:55.}

\teachervoice{El ponente enfatiza que esto no es como el Necker cube: no es que creas que cambió la depth del mundo exterior, sino que ves el mismo hecho con otra estructura. Usa las dos lecturas sintácticas de «next weekend we should be visiting relatives» para mostrar que un representation change no necesariamente modifica la truth condition.}

\subsection{La intrinsic relation va en los weights; la viewer pose, en la activity}

«La relación fija entre la crown y los flaps» debe mantenerse invariante con el viewpoint, por lo que conviene representarla con learned weights; «la pose de la crown respecto del observador», en cambio, varía con el punto de vista, por lo que conviene representarla con dynamic activity. Las $R_{wx}$, $R_{wv}$ y $R_{xv}$ de la figura de la clase pueden entenderse como transformaciones entre los object/part/viewer frames.

\lecturefigure{slide-12-crown-parse-tree.jpg}{Crown parse tree: las aristas guardan la parent-child intrinsic coordinate relation.}{Video oficial de Stanford Online 00:17:56--00:17:59.}

\lecturefigure{slide-13-zigzag-parse-tree.jpg}{Zig-zag/rectangle parse: los mismos rods corresponden a otro árbol estructural.}{Video oficial de Stanford Online 00:18:00--00:20:41.}

\begin{knowledgebox}{Lectura de la figura: comparar el node grouping, no la posición de los segmentos}
Los nodos hoja de ambos árboles siguen siendo el mismo conjunto de seis rods; lo que cambia son los nodes intermedios: el crown parse organiza los tres flaps dentro de un mismo whole, mientras que el zig-zag parse destaca el rectangle y los lóbulos inclinados superior e inferior. La $R$ de las aristas representa la intrinsic frame relation. Al leer la figura, primero se compara qué hojas comparten parent y después qué relations deben mantenerse invariantes con el viewpoint; no hay que confundir las dos figuras con dos objetos externos distintos.
\end{knowledgebox}

Una composition simplificada puede escribirse así:

\[
R_{xv}=R_{wv}R_{xw}.
\]

Aquí, $R_{xw}$ representa la intrinsic relation de la part $x$ respecto del whole $w$, $R_{wv}$ representa la pose actual del whole respecto del viewer $v$ y $R_{xv}$ es la part-to-viewer pose que se deduce. Distintas conventions pueden cambiar el orden de la multiplicación; lo que importa didácticamente es separar el viewpoint-invariant knowledge del viewpoint-dependent state.

\begin{warningbox}{La matrix equation no es un compromiso de implementación de la clase}
El artículo de GLOM analiza varias representaciones de pose/identity y no exige que toda relation sea una matriz explícita de transformación rígida. El embedding real también debe dar cabida a identity, uncertainty, texture y estructuras no rígidas. Forzar toda la semántica dentro de $SE(3)$ sería una simplificación excesiva.
\end{warningbox}

\subsection{La mental image no es un pixel buffer}

La postura de Hinton sobre la mental imagery es que una imagen interna puede ser una structural description más un conjunto de viewer-relative transforms. Para responder preguntas espaciales, las personas eligen implícitamente orientation, scale y position; si se cambia de frame, también cambian la ruta de razonamiento y las «relaciones evidentes».

\lecturefigure{slide-14-mental-image-crown.jpg}{Mental image: se añade la node-to-viewer pose al árbol estructural para facilitar la propagación del viewpoint.}{Video oficial de Stanford Online 00:20:42--00:22:01.}

\teachervoice{En clase se pide al público imaginar «una milla hacia el este, una milla hacia el norte y otra milla hacia el este». La mayoría supone por defecto que el north apunta hacia arriba y elige cierta escala y cierta posición en la imagen, aunque resolver el problema no exige ninguna de esas elecciones. Este pequeño experimento revela que el reference frame suele ser una variable latente del reasoning.}

\subsection{Resumen de la sección}

La part-whole hierarchy necesita expresar a la vez identity y pose. GLOM busca colocar el relation knowledge estable en la función de la red, y el viewpoint y el parse específicos de cada entrada en la embedding activity; el cube, los six rods y la mental imagery aportan una motivación cognitiva para esta división de funciones, pero no sustituyen los experimentos con el modelo.
